\documentclass{report}

% Page Size
\usepackage[letterpaper, portrait, margin=1in, left=1.5in]{geometry}

% Spacing
\usepackage{setspace}
\setlength{\parskip}{\baselineskip}
\doublespacing{}


% Bibliography
\usepackage[american]{babel}
\usepackage[style=apa, citestyle=apa, backend=biber]{biblatex}
\DeclareLanguageMapping{american}{american-apa}
\addbibresource{references.bib}
\usepackage[babel,threshold=2]{csquotes}

% Graphics
\usepackage{graphicx}
\graphicspath{ {./images/} }

% Table of Content
\setcounter{tocdepth}{2}

% Code Formatting
\usepackage{minted}
\usemintedstyle{pastie}

\def\code#1{\texttt{#1}}

% Positioning
\usepackage{float}

% Colors
\usepackage[dvipsnames]{xcolor}

% Links
\usepackage{hyperref}
\hypersetup{
	pdftitle={Development of a touch typing trainer with an emphasis on finger and wrist placement},
	pdfstartview={FitH},
	colorlinks=true,
	linkcolor=black,
	citecolor=black,
	filecolor=black,
	urlcolor=MidnightBlue
}


\begin{document}

\pagenumbering{gobble}
\begin{titlepage}
	\centering
	\includegraphics[width=0.25\textwidth]{upc.png}
	\includegraphics[width=0.25\textwidth]{dcs.png}
	\par\vspace{1cm}
	{\scshape\LARGE Research Brief \par}
	\par\vspace{1cm}
	\begin{tabular}{ | l | l | }
		\hline
		Name               & Oscar Vian L. Valles                  \\
		\hline
		Student ID number  & 2018--06491                           \\
		\hline
		Project Title      & Development of a touch typing trainer \\ & with an emphasis on finger and wrist positions\\
		\hline
		Project No.        & 9                                     \\
		\hline
		Name of Supervisor & Dhong Fhel K. Gom-os                  \\
		\hline
	\end{tabular}




\end{titlepage}
\newpage

\tableofcontents
\newpage

\pagenumbering{arabic}

\setcounter{chapter}{1}
\section{Background of the Study}
There are a lot of educational typing tests available that help people learn
touch typing, including Monkeytype, TypeRacer, or Keybr. These typing tests list
out words that are then typed out. The inputted keys are then compared to check
if the user has typed the expected letter. At the end of the test, the time
taken is calculated and certain metrics are given. These metrics include words
per minute (WPM) and accuracy \parencite{bartnik2021}.

However, this method of examination leaves out a crucial part of typing —
ergonomics. Ergonomic typing prevents a lot of health issues in the future like
repetitive strain injury or carpal tunnel. One important factor that affects
ergonomics is the typing procedure and posture. This means the proper placement of
the wrist, hands, and hitting the keys using the right finger that is assigned
to the key.

Correct finger placement is usually taught at the beginning using a diagram,
with each key being associated with a specific finger. For instance, the letter
Q in a QWERTY layout should be hit using the fifth digit of the left hand, and
this is shown by coloring the fifth digit and the key Q with the same color or
by placing the letters directly on the fingers \parencite{dobson2009touch}.

Incorrect finger placement may cause these hand and wrist positions: ulnar
deviation, forearm pronation, and wrist extension \parencite{serina1999}. These
three are hand and wrist positions that are common in all activities, however,
prolonged periods in these positions may cause injuries such as Carpal
tunnel syndrome (CTS) \parencite{toosi2015}

In addition, this type of typing is frequently taught in the beginner level
\parencite{donica2018}. This means that there is a need to weed out bad habits
that may develop, like using the index finger for pressing the spacebar or
backspace. However, it is impractical for an educator to check each student if
they are not performing these movements as these may only show for a small
period which may not be caught in time.

Thus, there is a need for automatically detecting which finger is used during
typing, and for the position of the wrist in relation to the arm. One way to do
this is through finger and hand tracking. One solution for tracking is by using
image processing and machine learning. An example of this is MediaPipe
by~\cite{mediapipe}.

MediaPipe allows for various applications for machine learning in the field of
image processing. This includes, hand tracking, pose estimation, object
detection, and others. Another example of a library that allows for hand and
finger tracking is OpenCV by~\cite{opencv}. This is a tool that simplifies
computer vision and image processing. Machine learning can also be used with
OpenCV.

\section{Research Objectives}

\subsection{General Objectives}
To create a touch typing trainer that corrects poor finger placement and hand
position to develop better typing habits and healthier typing ergonomics.

\subsection{Specific Objectives}
\begin{itemize}
	\item To develop a program that tracks fingers and hand positions while typing
	\item To create a subroutine that ascertains which finger was used to type a key
	\item To detect if incorrect fingers were used to press a key or if the
	      position of the hand in relation to the wrist is problematic
	\item To assess the accuracy and performance of the developed program's ability
	      to perform the previously stated objectives
	\item To develop a user-friendly interface for users to train touch typing
	      using the developed program
	\item To generate key statistics of a user's touch typing performance
\end{itemize}


\section{Scope and Limitation}
This research will focus on typing on a 60\% keyboard. This type of keyboard
only has the alphanumeric part of the keyboard. This limits the number of keys
to be checked and the expected movement of the hand. Furthermore, The keycaps
will also be of a light color, while the surface that the keyboard rests upon
will be of a dark color.

In addition, the keyboard layout will be ANSI. This layout is described by the
American National Standards Institute \parencite{ansi}. This is the most
common layout in the United States. However, it is also used in numerous
English-speaking countries such as the Philippines, Malaysia, and India.

The program will expect the that user has all ten digits and has no hand,
finger, or wrist deformities. In addition, only the placement of the hands,
fingers, and wrists will be taken into account when determining if the
ergonomics of the user while typing is healthy. The program will not check
seating position, angle of elbows, and other metrics for an ergonomic typing
posture while typing.

Capturing of the video to be analyzed by the program would be limited to a
single 1080p webcam that is capturing in 60 frames per second. The camera will
be pointed downwards facing the keyboard and the hand. This means that the
vertical angle of the wrist may not be accurate.

\section{Significance of the Research}
This research is beneficial for all users of physical keyboards. These include a
vast majority of the population as there are a lot of professions that heavily
rely on keyboards. Examples include developers, physicians, educators,
accountants. By having better ergonomics while typing, wrist injuries can be
prevented, and typing speed may be increased

This research also helps educators, especially early educators teaching beginner
typers. By automatically checking for ergonomics, posture, and correct
technique, the burden of checking each student is lessened, and directed
interventions for bad habits can be easily created as students with these bad
habits are easily identified

This research has a direct impact on people that has hand or wrist injuries that
are caused by poor typing habits. By correcting these poor habits, pain from
these injuries will be lessened, and even be prevented from occurring in the
first place. A specific example of this is by reducing ulnar deviation which
affects the nerve that is indicative of CTS \parencite{toosi2015}.

\newpage
\nocite{*} \printbibliography[heading=bibintoc,title={References}]{}

\end{document}
